% ===== §25 =====
\section{The Hermeneutics of Private Culture}
\label{sec:hermeneutics}

\noindent
The mechanisms of intimate culture have been set out; the phenomenology closes with the question of how the meaning of what they generate is interpreted. This is the hermeneutic layer, and it is not a mere appendix to the account of mechanisms, for it discloses something the mechanisms alone do not: that the meaning of an intimate culture is not a content its forms contain but an activity its participants continually perform, and that this fact has consequences, ethical as much as epistemological, which the final part will take up. The section describes the interpretive life of intimate culture, the two ways it fails, and the recognition that its right interpretation is.

\subsection*{Meaning as an activity, not a content}

The private forms of an intimate culture do not contain their meaning the way a vessel contains a liquid, such that the meaning would be there in the form to be extracted by anyone who opened it. Their meaning is generated, and continually regenerated, in the ongoing interpretive activity of the two who share them---in the living use, the re-enactment, the mutual reading by which a private word is again meant and again understood. A private idiom is meaningful only so long as it is actively meant and taken; the meaning is not stored in it between uses, waiting, but is remade each time the two enact the little circuit of address and understanding that the idiom serves. This is why the interpretation of intimate culture is not, in the end, separable from its generation: to interpret a private form rightly is to participate in the activity that keeps it meaningful, and the interpretation is itself a continuation of the generating. It follows that intimate culture is meaning only in the present tense of its interpretation. When the interpretive activity stops---when the two no longer actively mean and take their shared forms---the meaning does not persist in the forms as a dormant content that might later be revived; it ceases, because it never was a content but only ever an activity, and an activity that has stopped is simply not occurring. This is the hermeneutic face of the fragility the previous pillar described: the culture dies when its interpretation stops, because its life \emph{was} the interpreting.

\subsection*{The two misrecognitions}

Against this account of meaning as living activity, two characteristic failures can be defined, and both are forms of \emph{misrecognition}---the mistaking of a living, generated process for a finished thing. The first is \emph{reification} in the direction of possession: the treatment of the shared culture as an object one \emph{has}, a possession to be held and surveyed, rather than a process one \emph{does} with the other. When a partner comes to regard the private culture as a thing they own---an inventory of idioms and rites that exists, complete, to be pointed at---the living interpretive activity has been replaced by the contemplation of a supposed object, and the culture begins to die in that partner's hands, for it lived only in the doing that the having has suspended. This is the failure the epistemological section already glimpsed, in which participatory knowing lapses into objectifying knowledge; here it is named as an interpretive failure with an ethical edge, for it is a way of ceasing to participate while appearing to honor the culture, a withdrawal disguised as possession.

The second misrecognition runs in the direction of exhibition, and it points directly toward the capture the final part will analyze: the treatment of the relation's internal meaning as a \emph{capital} to be displayed and converted, an asset to be shown to others for recognition or advantage. When the private forms of a bond are taken outside the circuit that charges them and put on display---performed for an audience, exhibited as evidence of the relationship's quality, converted into a currency of social standing---they have been misrecognized as the very thing the paper argued they cannot be: transferable, extractable, a capital. This misrecognition is not merely a category error; it is the point of entry for the capture of intimate culture by forces outside the couple, for it is precisely by inducing the partners to treat their bond's meaning as an exhibitable and convertible asset that the outer order draws that meaning out of the circuit where it lived and into a market where it can be valued and extracted. The two misrecognitions are thus the hermeneutic form of the two dangers the paper tracks: reification toward possession is the death of the culture from within, and reification toward exhibition is the door through which it is captured from without.

\subsection*{Interpretation as recognition}

Set against these failures, the right interpretation of an intimate culture has, in the Hegelian line, a determinate character: it is a mutual recognition. To interpret the shared culture rightly is not to decode a content but to participate in the activity in which each, reading the shared form, recognizes again who they are to the other and who the other is to them. The private idiom rightly interpreted is a mirror in which the relation knows itself: in meaning and taking the shared word, each affirms the bond the word arose from and re-cognizes the other as the one with whom the word is shared. Interpretation, in this sense, is the ongoing self-knowledge of the relation carried on through its exteriorized forms, the moment in the Hegelian circuit at which the relation, having gone out of itself into a culture, returns to itself in the reading of that culture and knows itself there. And the two misrecognitions are, correspondingly, the stalling of this movement: reification toward possession stops the return by freezing the exteriorized culture into a dead object that no longer reflects the relation back to itself, and reification toward exhibition diverts the return outward, into a display for others, so that the relation no longer knows itself in its culture but seeks to be known by others through it. Right interpretation keeps the circuit turning; misrecognition, in either direction, stops it. With this the interpretive life of intimate culture is described, and the phenomenology can take its final step, to the edge where the couple's culture meets, and must negotiate with, the outer symbolic order---the edge at which the second misrecognition becomes the doorway to capture, and at which the particular returns to the general the final part requires.
